\section{Applications} 
\label{Sec5:Application}
% Practical applications (We can discuss each application from Agent specific Features, Agent organization, environment, ……. etc perspectives. Some applications may only emphasize some perspectives (not all perspectives). For example, there is no environment in multi-agent software development).  
% Overview Table (Taicheng)
% Please add the following required packages to your document preamble:
% \usepackage{multirow}

% Please add the following required packages to your document preamble:
% \usepackage{multirow}
% \usepackage[table,xcdraw]{xcolor}
% Beamer presentation requires \usepackage{colortbl} instead of \usepackage[table,xcdraw]{xcolor}
% Please add the following required packages to your document preamble:
% \usepackage{multirow}
% \usepackage[table,xcdraw]{xcolor}
% Beamer presentation requires \usepackage{colortbl} instead of \usepackage[table,xcdraw]{xcolor}

LLM-MA systems have been used in a wide range of applications. We summarize two kinds of applications in Table \ref{summary}: \textbf{Problem Solving} and  \textbf{World Simulation}. We   elaborate on these applications below. Note that this is a fast growing research field and new applications appear almost everyday. We maintain an  \href{https://github.com/taichengguo/LLM_MultiAgents_Survey_Papers}{open source repository} to report the latest work.  

\subsection{LLM-MA for Problem Solving}

% What is for problem solving and some simple examples
%One main stream application of LLM-MA systems is in specific problem-solving scenarios, such as Software Development, Embodied Agents, Science Experiments and Science Debate.
The main motivation of using LLM-MA for problem solving is to harness the collective capabilities of agents with specialized expertise. These agents, each acting as individuals, collaborate to address complex problems effectively,  such as software development, embodied agents, science experiments and science debate. These application examples are introduced next.
% General Methodology: How to achieve this from four taxonomy perspective
% In the context of LLM-MA for problem solving, the environment is typically Sandbox or Physical . Each agent assumes a distinct, specialized role, working cooperatively to solve the target task.
% More details as follows
%Subsequent sections delve into a detailed explanation of LLM-MA for problem solving applications.


\subsubsection{Software Development}
% Summary: 1.Why this problem needs LLM-MA  2. How LLM-MA do in this case (Schema from section 3)
Given that software development is a complex endeavor requiring the collaboration of various roles like product managers, programmers, and testers, LLM-MA systems are typically set to emulate these distinct roles and collaborate to address the intricate challenge. Following the waterfall or Standardized Operating Procedures (SOPs) workflow of the software development, the communication structure among agents is usually layered. Agents generally interact with the code interpreter, other agents or human to iteratively refine the generated code.
% Each work
\cite{li2023camel} first proposes a simple role-play agent framework, which utilizes the interplay of two roles to realize autonomous programming based on one-sentence user instruction. It provides insights into the ``cognitive" processes of communicative agents. \cite{dong2023selfcollaboration} makes LLMs work as distinct ``experts" for sub-tasks in software development, autonomously collaborating to generate code. Moreover, \cite{qian2023communicative} presents an end-to-end framework for software development, utilizing multiple agents for software development without incorporating advanced human teamwork experience. \cite{hong2023metagpt} first incorporates human workflow insights for more controlled and validated performance. It encodes SOPs into prompts to enhance structured coordination. \cite{huang2023agentcoder} delves deeper into multi-agent based programming by solving the problem of balancing code snippet generation with effective test case generation, execution, and optimization. 

% Some key limitations for LLM-MA for Software Development discussed in previous work are:
% There are some general limitations and challenges. Notably, observed limitations include inherent randomness in generated output, resulting in variability across software productions in different runs. Furthermore, the LLMs may manifest biases, deviating from problem-solving approaches typical of real programmers. This bias introduces code patterns not necessarily aligned with authentic programming paradigms. These challenges are further compounded by the hallucinatory tendencies inherent in large language models, underscoring the complexities associated with ensuring accurate and unbiased code generation in the realm of automatic programming.

\subsubsection{Embodied Agents}
% Summary: 1.Why this problem needs LLM-MA  2. How LLM-MA do in this case (Schema from section 3)
Most embodied agents applications inherently utilize multiple robots working together to perform complex real-world planning and manipulation tasks such as warehouse management with heterogeneous robot capabilities. Hence, LLM-MA can be used to model robots with different capabilities and cooperate with each other to solve real-world physical tasks.
% Each Work
\cite{dasgupta2023collaborating} first explores the potential to use LLM as an action planner for embedded agents. \cite{mandi2023roco} introduces RoCo, a novel approach for multi-robot collaboration that uses LLMs for high-level communication and low-level path planning. Each robotic arm is equipped with an LLM, cooperating with inverse kinematics and collision checking. Experimental results demonstrate the adaptability and success of RoCo in collaborative tasks. \cite{zhang2023building} presents CoELA, a Cooperative Embodied Language Agent, managing discussions and task planning in an LLM-MA setting. This challenging setting is featured with decentralized control, complex partial observation, costly communication, and multi-objective long-horizon tasks. \cite{chen2023scalable} investigates communication challenges in scenarios involving a large number of robots, as assigning each robot an LLM will be costly and unpractical due to the long context. The study compares four communication frameworks, centralized, decentralized, and two hybrid models, to evaluate their effectiveness in coordinating complex multi-agent tasks. \cite{yu2023conavgpt} proposes Co-NavGPT for multi-robot cooperative visual target navigation, integrating LLM as a global planner to assign frontier goals to each robot. \cite{chen2023multi} proposes an LLM-based consensus-seeking framework, which can be applied as a cooperative planner to a multi-robot aggregation task.


% The limitations 
% There are some general limitations. Current works for LLMs-based Embodied Agents mainly focus on high-level task planning, considering the inefficiency and poor performance of LLMs for low-level actions \cite{zhang2023building}. Besides, it is usually needed that the environment is already grounded in language \cite{nottingham2023embodied}, and multi-modal large models capable of both processing visual modalities effectively and generating natural language fluently \cite{huang2023language}\cite{driess2023palm} \cite{lu2022unified} would help. It is also mentioned that multiple rounds of discussions in a multi-agent system are inefficient where they might have difficulty reaching a consensus or repeat historical ideas and wrong actions \cite{chen2023scalable}.

\subsubsection{Science Experiments}
% ChatGPT team; Medical Team; 
% Motivation: Why use LLM-based Multi-Agents to solve this problem
Like multiple agents play as different specialists and cooperate to solve the Software Development and Embodied Agents problem, multiple agents can also be used to form a science team to conduct science experiments. One important difference from previous applications lies in the crucial role of human oversight,  due to the high expenses of the science experiments and the hallucination of the LLM agents. Human experts are at the center of these agents to process the information of agents and give feedback to the agents.
% How they use Multi-Agents; Conclusion
~\cite{chatgptforchem} utilizes multiple LLM-based agents, each focusing on specific tasks for the science experiments including strategy planning, literature search, coding, robotic operations, and labware design. All these agents interact with humans to work collaboratively to optimize the synthesis process of complex materials.

\subsubsection{Science Debate}
LLM-MA can be set for science debating scenarios, where agents debate with each other to enhance the collective reasoning capabilities in tasks such as Massive Multitask Language Understanding (MMLU)~\cite{hendrycks2020measuring}, Math problems~\cite{cobbe2021training}, and StrategyQA~\cite{geva2021did}. The main idea is that each agent initially offers its own analysis of a problem, which is then followed by a joint debating process. Through multiple rounds of debate, the agents converge on a single, consensus answer. 
~\cite{du2023improving} leverages the multi-agents debate process on a set of six different reasoning and factual accuracy tasks and demonstrates that LLM-MA debating can improve factuality. ~\cite{xiong2023examining} focuses on the commonsense reasoning tasks and formulates a three-stage debate to align with real-world scenarios including fair debate, mismatched debate, and roundtable debate. The paper also analyzes the inter-consistency between different LLMs and claims that debating can improve the inter-consistency.~\cite{tang2023medagents} also utilizes multiple LLM-based agents as distinct domain experts to do the collaborative discussion on a medical report to reach a consensus for medical diagnosis. 


\subsection{LLM-MA for World Simulation}

% What is for problem solving and some simple examples
Another mainstream application scenario of LLM-MA is the world simulation.  Research in this area is rapidly growing and spans a diverse range of fields including social sciences, gaming, psychology, economics, policy-making, etc.
% Why we need Multi-Agents in this case: LLM strong ability
The key reason for employing LLM-MA in world simulations lies in their exceptional role-playing abilities, which are crucial for realistically depicting various roles and viewpoints in a simulated world.
% General Methodology: How to achieve this from four taxonomy perspective
The environment of world simulation projects is usually crafted to reflect the specific scenario being simulated, with agents designed in various profiles to match this context. Unlike the problem solving systems that focus on agent cooperation, world simulation systems involve diverse methods of agent management and communication, reflecting the complexity and variety of real-world interactions.
% More details as follows
Next, we explore simulations conducted in diverse fields.
%The following sections provide in-depth information and analysis.



\subsubsection{Societal Simulation}
In societal simulation, LLM-MA models are used to simulate social behaviors, aiming to explore the potential social dynamics and propagation, test social science theories, and populate virtual spaces and communities with realistic social phenomena~\cite{park2023generative}.  Leveraging LLM's capabilities, agents with unique profiles engage in extensive communication, generating rich behavioral data for in-depth social science analysis. 

% Current Work
The scale of societal simulation has expanded over time, beginning with smaller, more intimate settings and progressing to larger, more intricate ones. 
Initial work by \cite{park2023generative} introduces generative agents within an interactive sandbox environment reminiscent of the sims, allowing end users to engage with a modest community of 25 agents through natural language. 
At the same time, \cite{park2022social} develops Social Simulacra, which constructs a simulated community of 1,000 personas. This system takes a designer's vision for a community—its goals, rules, and member personas—and simulates it, generating behaviors like posting, replying, and even anti-social actions.
Building on this, \cite{gao2023s} takes the concept further by constructing vast networks comprising 8,563 and 17,945 agents, respectively, designed to simulate social networks focused on the topics of Gender Discrimination and Nuclear Energy. This evolution showcases the increasing complexity and size of simulated environments in recent research.
Recent studies such as \cite{chen2023multi,kaiya2023lyfe,li2023quantifying,li2023you,ziems2023can} highlight the evolving complexity in multi-agent systems, LLM impacts on social networks, and their integration into social science research.
 
% Emotion 
% Information
\subsubsection{Gaming}
% Summary: 1. What is Game theory 2. Why this problem needs LLM-MA  3. How LLM-MA normally do in this case (Schema from section 3)
LLM-MA is well-suited for creating simulated gaming environments, allowing agents to assume various roles within games. This technology enables the development of controlled, scalable, and dynamic settings that closely mimic human interactions, making it ideal for testing a range of game theory hypotheses~\cite{mao2023alympics,xu2023exploring,gong2023mindagent}. 
%The Game theory is the study of models of strategic interaction between rational decision-makers~\cite{game_theory}. LLM-MA is suitable for game theory simulation since it can facilitate the construction of a controlled, scalable, realistic and dynamic environment of human interactions for many game theory hypothesis testing~\cite{mao2023alympics,xu2023exploring}. 
Most games simulated by LLM-MA rely heavily on natural language communication, offering a sandbox environment within different game settings for exploring or testing game theory hypotheses including reasoning, cooperation, persuasion, deception, leadership, etc. 

~\cite{akata2023playing} leverages behavioral game theory to examine LLMs' behavior in interactive social settings, particularly their performance in games like the iterated Prisoner's Dilemma and Battle of the Sexes. Furthermore, \cite{xu2023exploring} proposes a framework using ChatArena library \cite{ChatArena} for engaging LLMs in communication games like Werewolf, using retrieval and reflection on past communications for improvement, as well as the Chain-of-Thought mechanism \cite{wei2022chain}.
~\cite{light2023text} explores the potential of LLM agents in playing Resistance Avalon, introducing AVALONBENCH, a comprehensive game environment and benchmark for further developing advanced LLMs and multi-agent frameworks. ~\cite{wang2023avalon} also focuses on the capabilities of LLM Agents in dealing with misinformation in the Avalon game, proposing the Recursive Contemplation (ReCon) framework to enhance LLMs' ability to discern and counteract deceptive information. 
~\cite{xu2023language} introduces a framework combining LLMs with reinforcement learning (RL) to develop strategic language agents for the Werewolf game. It introduces a new approach to use RL policy in the case that the action and state sets are not predefined but in the natural language setting. ~\cite{mukobi2023welfare} designs the ``Welfare Diplomacy”, a general-sum variant of the zero-sum board game Diplomacy, where players must balance military conquest and domestic welfare. It also offers an open-source benchmark, aiming to help improve the cooperation ability of multi-agent AI systems. On top of that, there is a work \cite{li2023theory} in a multi-agent cooperative text game testing the agents' Theory of Mind (ToM), the ability to reason about the concealed mental states of others and is fundamental to human social interactions, collaborations, and communications. \cite{fan2023can} comprehensively assesses the capability of LLMs as rational players, and identifies the weaknesses of LLM-based Agents that even in the explicit game process, agents may still overlook or modify refined beliefs when taking actions.


\subsubsection{Psychology}
In psychological simulation studies, like in the societal simulation, multiple agents are utilized to simulate humans with various traits and thought processes. However, unlike societal simulations, one approach in psychology involves directly applying psychological experiments to these agents. This method focuses on observing and analyzing their varied behaviors through statistical methods. Here, each agent operates independently, without interacting with others, essentially representing different individuals. Another approach aligns more closely with societal simulations, where multiple agents interact and communicate with each other. In this scenario, psychological theories are applied to understand and analyze the emergent behavioral patterns. This method facilitates the study of interpersonal dynamics and group behaviors, providing insights into how individual psychological traits influence collective actions.
~\cite{ma2023understanding} explores the psychological implications and outcomes of employing LLM-based conversational agents for mental well-being support. It emphasizes the need for carefully evaluating the use of LLM-based agents in mental health applications from a psychological perspective. 
~\cite{kovavc2023socialai} introduces a tool named SocialAI school for creating interactive environments simulating social interactions. It draws from developmental psychology to understand how agents can acquire, demonstrate, and evolve social skills such as joint attention, communication, and cultural learning. 
~\cite{zhang2023exploring} explores how LLM agents, with distinct traits and thinking patterns, emulate human-like social behaviors such as conformity and majority rule. This integration of psychology into the understanding of agent collaboration offers a novel lens for examining and enhancing the mechanisms behind LLM-based multi-agents systems.
~\cite{aher2023using} introduces Turing Experiments to evaluate the extent to which large language models can simulate different aspects of human behaviors. The Turing Experiments replicate classical experiments and phenomena in psychology, economics, and sociology using a question-answering format to mimic experimental conditions. They also design a prompt that is used to simulate the responses of multiple
different individuals by varying the name. By simulating various kinds of individuals via LLM, they show that larger models replicate human behavior more faithfully, but they also reveal a hyper-accuracy distortion, especially in knowledge-based tasks. 


% Please add the following required packages to your document preamble:
% \usepackage{multirow}
% \usepackage[normalem]{ulem}
% \useunder{\uline}{\ul}{}
\begin{table*}[ht]
\centering
\resizebox{0.9\textwidth}{!}{
\begin{tabular}{l|l|l|l|l}
\hline
\textbf{Motivation} &
  \textbf{Domain} &
  \textbf{Datasets and Benchmarks} &
  \textbf{Used by} &
  \textbf{Data Link} \\ \hline
\multirow{13}{*}{Problem Solving} &
  \multirow{3}{*}{Software Development} &
  HumanEval &
  ~\cite{hong2023metagpt} & \href{https://github.com/openai/human-eval}{Link} \\
 &
   &
  MBPP &
  ~\cite{hong2023metagpt} & \href{https://github.com/google-research/google-research/tree/master/mbpp}{Link}
   \\
 &
   &
  SoftwareDev &
  ~\cite{hong2023metagpt} &
  \href{https://github.com/geekan/MetaGPT}{Link} \\ \cline{2-5} 
 &
  \multirow{4}{*}{Embodied AI} &
  RoCoBench &
  ~\cite{mandi2023roco} &
  \href{https://github.com/MandiZhao/robot-collab/tree/main/rocobench}{Link} \\
 &
   &
  Communicative Watch-And-Help (C-WAH) &
  ~\cite{zhang2023building} &
  \href{https://github.com/UMass-Foundation-Model/Co-LLM-Agents/}{Link} \\
 &
   &
  ThreeDWorld Multi-Agent Transport (TDW-MAT) &
  ~\cite{zhang2023building} &
  \href{https://github.com/UMass-Foundation-Model/Co-LLM-Agents/}{Link} \\
 &
   &
  HM3D v0.2 &
  ~\cite{yu2023conavgpt} &
  \href{https://aihabitat.org/datasets/hm3d-semantics/}{Link} \\ \cline{2-5} 
 &
  \multirow{6}{*}{Science Debate} &
  MMLU &
  ~\cite{tang2023medagents} &
  \href{https://github.com/hendrycks/test}{Link} \\
 &
   &
  MedQA &
   ~\cite{tang2023medagents} &
  \href{https://github.com/jind11/MedQA}{Link} \\
 &
   &
  PubMedQA &
  ~\cite{tang2023medagents} &
  \href{https://github.com/pubmedqa/pubmedqa}{Link} \\
 &
   &
  GSM8K &
  ~\cite{du2023improving} &
  \href{https://github.com/openai/grade-school-math}{Link} \\
 &
   &
  StrategyQA &
  ~\cite{xiong2023examining} &
  \href{https://github.com/eladsegal/strategyqa}{Link} \\
 &
   &
  Chess Move Validity &
  ~\cite{du2023improving} &
  \href{https://github.com/google/BIG-bench}{Link} \\ \hline
\multirow{15}{*}{World Simulation} &
  \multirow{3}{*}{Society} &
  SOTOPIA &
  ~\cite{zhou2023sotopia} &
  / \\
 &
   &
  Gender Discrimination &
  ~\cite{gao2023s} &
  / \\
 &
   &
  Nuclear Energy &
  ~\cite{gao2023s} &
  / \\ \cline{2-5} 
 &
  \multirow{6}{*}{Gaming} &
  Werewolf &
  ~\cite{xu2023exploring} &
  / \\
 &
   &
  Avalon &
  ~\cite{light2023text} &
  / \\
 &
   &
  Welfare Diplomacy &
  ~\cite{mukobi2023welfare} &
  / \\
 &
   &
  Layout  in the Overcooked-AI environment &
  ~\cite{agashe2023evaluating} &
  / \\
 &
   &
  Chameleon &
  ~\cite{xu2023magic} &
  \href{https://github.com/cathyxl/MAgIC}{Link} \\
 &
   &
  Undercover &
  ~\cite{xu2023magic} &
  \href{https://github.com/cathyxl/MAgIC}{Link} \\ \cline{2-5} 
 &
  \multirow{3}{*}{Psychology} &
  Ultimatum Game TE &
  ~\cite{aher2023using} &
  \href{https://github.com/microsoft/turing-experiments/tree/main}{Link} \\
 &
   &
  Garden Path TE &
  ~\cite{aher2023using} &
  \href{https://github.com/microsoft/turing-experiments/tree/main}{Link} \\
 &
   &
  Wisdom of Crowds TE &
  ~\cite{aher2023using} &
  \href{https://github.com/microsoft/turing-experiments/tree/main}{Link} \\ \cline{2-5} 
 &
  \multirow{2}{*}{Recommender System} &
  MovieLens-1M &
  ~\cite{zhang2023generative} &
  \href{https://github.com/LehengTHU/Agent4Rec/tree/master/datasets/ml-1m}{Link} \\
 &
   &
  Amazon review dataset &
  ~\cite{zhang2023agentcf} &
  / \\ \cline{2-5} 
 &
  Policy Making &
  Board Connectivity Evaluation &
  ~\cite{hua2023war} &
  \href{https://github.com/agiresearch/WarAgent}{Link} \\ \hline
\end{tabular}
}
\caption{Datasets and Benchmarks commonly used in LLM-MA studies. `` / " denotes the unavailability of  data link.}
\label{db}
\end{table*}


\subsubsection{Economy}
LLM-MA is used to simulate economic and financial trading environments mainly because it can serve as implicit computational models of humans. In these simulations, agents are provided with endowments, and information, and set with pre-defined preferences, allowing for an exploration of their actions in economic and financial contexts. This is similar to the way economists model 'homo economicus', the characterization of man in some economic theories as a rational person who pursues wealth for his own self-interest ~\cite{horton2023large}. 
There are several studies demonstrate the diverse applications of LLM-MA in simulating economic scenarios, encompassing macroeconomic activities, information marketplaces, financial trading, and virtual town simulations. Agents interact in cooperative or debate, decentralized environments. \cite{li2023large} employs LLMs for macroeconomic simulation, featuring prompt-engineering-driven agents that emulate human-like decision-making, thereby enhancing the realism of economic simulations compared to rule-based or other AI agents. \cite{anonymous2023rethinking} explores the buyer's inspection paradox in an information marketplace, reveals improved decision-making and answer quality when agents temporarily access information before purchase. \cite{li2023tradinggpt} presents an LLM-MA framework for financial trading, emphasizing a layered memory system, debate mechanisms, and individualized trading characters, thereby fortifying decision-making robustness. \cite{zhao2023competeai} utilizes LLM-based Agents to simulate a virtual town with restaurant and customer agents, yielding insights aligned with sociological and economic theories. These studies collectively illuminate the broad spectrum of applications and advancements in employing LLMs for diverse economic simulation scenarios.


\subsubsection{Recommender Systems}
% Representative works
% Motivation; How they use Multi-Agents; 
The use of the LLM-MA in recommender systems is similar to that in psychology since studies in both fields involve the consideration of extrinsic and intrinsic human factors such as cognitive processes and personality~\cite{PsychologyInformedRecommenderSystems}. One way to use LLM-MA in recommender systems is to directly introduce items to multiple LLM-based agents within diverse traits and conduct statistics of the preferences of different agents. Another way is to treat both users and items as agents and the user-item communication as interactions, simulating the preference propagation. To bridge the gap between offline metrics and real-world performance in recommendation systems, Agent4Rec~\cite{zhang2023generative} introduces a simulation platform based on LLM-MA. 1000 generative agents are initialized with the MovieLens-1M dataset to simulate complex user interactions in a recommendation environment. 
% conclusion; 
Agent4Rec shows that LLM-MA can effectively mimic real user preferences and behaviors, provide insights into phenomena like the filter bubble effect, and help uncover causal relationships in recommendation tasks. In Agent4Rec work, agents are used to simulate users and they do not communicate with each other. Different from Agent4Rec work, ~\cite{zhang2023agentcf} treats both users and items as agents, optimizing them collectively to reflect and adjust to real-world interaction disparities. This work emphasizes simulating user-item interactions and propagates preferences among agents, capturing the collaborative filtering essence.

\subsubsection{Policy Making}
Similar to simulations in gaming and economic scenarios, Policy Making requires strong decision-making capabilities to realistic and dynamic complex problems. LLM-MA can be used to simulate the policy making via simulating a virtual government or simulating the impact of various policies on different communities. These simulations provide valuable insights into how policies are formulated and their potential effects, aiding policymakers in understanding and anticipating the consequences of their decisions~\cite{policy}.
The research outlined in  \cite{xiao2023simulating} is centered on simulating a township water pollution crisis. It simulated a town located on an island including a demographic structure of different agents and township head and advisor. Within the water pollution crisis simulation, this work provides an in-depth analysis of how a virtual government entity might respond to such a public administration challenge and how information transfer in the social network in this crisis. ~\cite{hua2023war} introduces WarAgent to simulate key historical conflicts and provides insights for conflict resolution and understanding, with potential applications in preventing future international conflicts.
% The use of agent simulations in policy-making acts as a valuable roadmap for developing effective diffusion models that incorporate human reasoning and decision-making processes.


\subsubsection{Disease Propagation Simulation}
Leveraging the societal simulation capabilities of LLM-MA can also be used to simulate disease propagation.
The most recent study in \cite{williams2023epidemic} delves into the use of LLM-MA in simulating disease spread. 
% In this research, agents are uniquely programmed to autonomously reason and make decisions, facilitated through their connection to ChatGPT. 
The research showcases through various simulations how these LLM-based agents can accurately emulate human responses to disease outbreaks, including behaviors like self-quarantine and isolation during heightened case numbers. 
The collective behavior of these agents mirrors the complex patterns of multiple waves typically seen in pandemics, eventually stabilizing into an endemic state. 
Impressively, their actions contribute to the attenuation of the epidemic curve. 
% This innovative approach marks a significant stride in dynamic system modeling, introducing a sophisticated representation of human cognitive processes, reasoning, and decision-making.
~\cite{ghaffarzadegan2023generative} also discusses the epidemic propagation simulation and decomposes the simulation into two parts: the Mechanistic Model which represents the information or propagation of the virus and the Decision-Making Model which represents the agents' decision-making process when facing the virus.

%\subsubsection{More Works..}
%Because the LLM-based Multi-Agents is a fast growing research field and new applications appear almost everyday, it is difficult to include each application in the review. Hence, we maintain an updated open source repository to collect the latest work. The repository link is \href{https://github.com/taichengguo/LLM_MultiAgents_Survey_Papers}{LLM\_MultiAgents\_Survey\_Papers}.

% \subsection{Multi-Agents for Dialogue Dataset Generation}





